\documentclass[11pt]{amsart}

\usepackage[margin=1.2in]{geometry}
\usepackage{amsmath,amssymb,amsthm}
\usepackage[colorlinks=true,linkcolor=black,citecolor=black,urlcolor=black]{hyperref}

\newtheorem{theorem}{Theorem}[section]
\newtheorem{lemma}[theorem]{Lemma}
\newtheorem{corollary}[theorem]{Corollary}
\newtheorem{proposition}[theorem]{Proposition}
\theoremstyle{remark}
\newtheorem{remark}[theorem]{Remark}
\theoremstyle{definition}
\newtheorem{definition}[theorem]{Definition}

\newcommand{\arxiv}[1]{\href{https://arxiv.org/abs/#1}{arXiv:#1}}
\newcommand{\N}{\mathbb N}
\newcommand{\C}{\mathbb C}
\newcommand{\RadMF}{\operatorname{Rad}_{\mathrm{MF}}}
\newcommand{\ncl}[1]{\langle\!\langle #1\rangle\!\rangle}

\title[Saturating the compression defect]{Saturating the compression defect:\\
a torsion-free finitely presented group\\ with no sofic and no MF quotients}
\author{Draft note for expert review --- not for citation}
\date{August 14, 2026}

\begin{document}

\begin{abstract}
We combine the construction of Fournier-Facio's torsion-free non-sofic
group with one further application of Hull's small-cancellation theorem for
acylindrically hyperbolic groups, routed so that the embedded simple factor
normally generates the resulting quotient.  The output is a
$2$-generated, finitely presented, torsion-free group $Q$ with property
\textup{(T)}, acylindrically hyperbolic, such that \emph{every} nontrivial
quotient of $Q$ --- including $Q$ itself --- is non-sofic and has full MF
radical: every homomorphism from a nontrivial quotient of $Q$ to the unitary
group of a norm matrix corona is trivial.  The two obstruction halves
(the norm-corona radical theorem and the nonsoficity criterion) and all
algebraic packaging are machine-checked in Lean; the geometric
small-cancellation existence input is cited from Hull, Osin, and
Fournier-Facio and is exactly the unformalized boundary.  This note isolates
that boundary for expert scrutiny.
\end{abstract}

\maketitle

\section{Status and purpose}\label{sec:status}

This note is a request for a sanity check, not an announcement.  It
assembles a strengthening of the construction in \cite{FFF} from three
components:

\begin{enumerate}
\item a machine-checked analytic theorem: a normal property-\textup{(T)}
  subgroup contained in the \emph{compression defect} of a group lies in its
  MF radical (\S\ref{sec:obstruction});
\item a machine-checked algebraic observation: in the Fournier-Facio
  configuration the compression defect is exactly the normal closure of the
  embedded simple factor (\S\ref{sec:ff});
\item a small-cancellation step, assembled entirely from published
  statements of Hull and Osin: the normal closure of the simple factor can be
  routed \emph{onto} a further quotient while preserving finite presentation,
  torsion-freeness, two loxodromic generators, and a protected finite set
  (\S\ref{sec:saturation}).
\end{enumerate}

Components (1) and (2), the nonsoficity criterion, and every implication
downstream of (3) are formalized in Lean~4 in the accompanying repository;
\S\ref{sec:boundary} states the exact certification boundary.  Component (3)
is the only part that is not formalized, and it is the part this note asks an
expert in acylindrical small cancellation to audit.  The main claim is
Theorem~\ref{thm:main}.

Throughout, \emph{MF group} has the Carri\'on--Dadarlat--Eckhardt meaning
\cite{CDE}: a countable discrete group $G$ is MF if it embeds in the unitary
group of a \emph{norm matrix corona}
$\mathcal Q=\prod_n M_{d_n}(\C)/\bigoplus_n M_{d_n}(\C)$
(bounded sequences modulo sequences vanishing in operator norm).  The
\emph{MF radical} $\RadMF(G)$ is the intersection of the kernels of all
homomorphisms $G\to\mathcal U(\mathcal Q)$ over all coronas; thus $G$ is
non-MF as soon as $\RadMF(G)\neq1$, and $\RadMF(G)=G$ says that every
homomorphism from $G$ to the unitary group of any norm matrix corona is
trivial.

\section{The machine-checked obstructions}\label{sec:obstruction}

\begin{definition}\label{def:core}
A \emph{compression core} on a countable group $G$ consists of a group
$\Gamma$ with property \textup{(T)}, a homomorphism
$\iota\colon\Gamma\to G$, an element $t\in G$ with
$t\,\iota(\Gamma)\,t^{-1}\subseteq\iota(\Gamma)$, and an element $c\in G$
commuting with $\iota(\Gamma)$.  Its \emph{defect} is the normal subgroup
\[
D \;=\; \ncl{\,[\,t c t^{-1},\ \iota(\gamma)\,]\ :\ \gamma\in\Gamma\,}_G .
\]
\end{definition}

If $c$ still commuted with $\iota(\Gamma)$ \emph{after} transport by $t$,
the defect would vanish; $D$ measures the failure of one-sided Kazhdan
transport to be implementable by an honest symmetry.

\begin{theorem}[machine-checked; \texttt{normalKazhdan\_le\_normMFResidual}]
\label{thm:radical}
Let $G$ be a countable group carrying a compression core with defect $D$,
and let $K\trianglelefteq G$ be a normal subgroup with property \textup{(T)}
such that $K\subseteq D$.  Then $K\subseteq\RadMF(G)$: every homomorphism
from $G$ to the unitary group of a norm matrix corona kills $K$ pointwise.
No finiteness, centrality, or torsion hypothesis on $K$ is used.
\end{theorem}

The second obstruction is the compression nonsoficity criterion, in the form
of \cite[Proposition~1.2]{FFF} (the same mechanism as the criterion driving
\cite{OAI,KT}): if $\Gamma\leq G$ are infinite groups with property
\textup{(T)}, $G=\langle\Gamma,t_1,\dots,t_m\rangle$ with
$t_i\Gamma t_i^{-1}\subseteq\Gamma$, and $J\leq G$ is finitely generated with
$[\Gamma,J]=1$, $\Gamma\cap J=1$, and $t_1Jt_1^{-1}\subseteq\Gamma$, then
soficity of $G$ forces $J$ to be LEF.  This criterion is also machine-checked
in the repository (\texttt{CriterionAssembly.not\_isSofic\_of\_not\_isLEF}
over the interface \texttt{CompressionSetup}).  Since a finitely presented
LEF group is residually finite (Vershik--Gordon) and an infinite simple
group is never residually finite, a finitely presented infinite simple $J$
is never LEF; this is the use made of it below, exactly as in \cite{FFF}.

\section{The Fournier-Facio configuration and its defect}\label{sec:ff}

We recall the output of \cite[\S2]{FFF}, in the notation of that paper.
Fournier-Facio constructs, from a universal finitely presented torsion-free
group $U$, a finitely presented infinite simple torsion-free group $S$
(Burger--Mozes or Hyde--Lodha), and a torsion-free hyperbolic Kazhdan group
$H$:

\begin{enumerate}
\item a finitely presented torsion-free property-\textup{(T)} group $P$
  containing $P_1\times P_2\times S$ with $P_i\cong P$, via small
  cancellation over the relatively hyperbolic pair $(U*H,\,U)$
  \cite[Theorem~2.4]{Osin10};
\item the double HNN extension $E=\langle P,u_1,u_2\mid u_iPu_i^{-1}=P_i\rangle$,
  torsion-free, finitely presented, and acylindrically hyperbolic by
  \cite{MO15};
\item a common quotient $\pi\colon E\twoheadrightarrow G$ of $E$ and $H$ via
  \cite[Corollary~7.4]{Hull16}, injective on a prescribed finite set
  containing some $s\in S\smallsetminus\{1\}$; $G$ is finitely presented,
  torsion-free by \cite[Theorem~7.1(e)]{Hull16}, and has property
  \textup{(T)} as a quotient of $H$.
\end{enumerate}

Write $\Gamma=\pi(P)$, $t_i=\pi(u_i)$, and $J=t_1^{-1}\pi(S)t_1$, so that
$[\Gamma,J]=1$, $\Gamma\cap J=1$, $t_1Jt_1^{-1}=\pi(S)\subseteq\Gamma$, and
$\pi|_S$ is injective by simplicity.  (These are the data entering
\cite[Proposition~1.2]{FFF}; the harmless discrepancy between
$t_1J t_1^{-1}$ and $t_1^{-1}Jt_1$ conventions is discussed in the companion
document cited in \S\ref{sec:boundary}.)

\begin{remark}[acylindrical hyperbolicity of $G$]\label{rem:ah}
The paper \cite{FFF} does not record whether $G$ is acylindrically
hyperbolic, because nothing there needs it.  We need it, and it is available
at no cost: since $E$ is finitely generated, the proof of
\cite[Corollary~7.4]{Hull16} produces the common quotient in the class
$\mathcal{AH}_0$ (two applications of \cite[Theorem~7.1]{Hull16}, the second
of which makes the image of the suitable subgroup all of $G$; Hull states
the $\mathcal{AH}_0$ conclusion explicitly in the finitely generated case).
We therefore fix once and for all a group $G$ produced by running
\cite[\S2]{FFF} with this clause of Hull's corollary invoked, so that
\[
G\ \text{is finitely presented, torsion-free, has property \textup{(T)},
and lies in } \mathcal{AH}_0 .
\]
This is a reading of the cited corollary, not a modification of the
construction; it is the first of the two points we would like an expert to
confirm.
\end{remark}

Now take the compression core on $G$ given by $\iota$ the inclusion of
$\Gamma$, $t=t_1$, and any $c\in J\smallsetminus\{1\}$, so
$d:=t_1ct_1^{-1}\in\pi(S)\smallsetminus\{1\}$.

\begin{lemma}[defect identification]\label{lem:defect}
$D=\ncl{\pi(S)}_G$.  The containment $\supseteq$ is machine-checked, as
\textup{\texttt{simpleSubgroup\_le\_defectNormal}}.
\end{lemma}

\begin{proof}
($\subseteq$) Each generator $[d,\iota(\gamma)]=d\,(\gamma d^{-1}\gamma^{-1})$
is a product of two conjugates of $d^{\pm1}\in\pi(S)$, hence lies in the
normal subgroup $\ncl{\pi(S)}_G$.

($\supseteq$) $D\cap\pi(S)$ is normal in $\pi(S)$.  Since $S$ is infinite
simple it is nonabelian with trivial centre, so $d\neq1$ is not central in
$\pi(S)$: there is $\gamma\in\pi(S)\subseteq\Gamma$ with $[d,\gamma]\neq1$.
That commutator lies in $D\cap\pi(S)$, which is therefore nontrivial, hence
equals $\pi(S)$ by simplicity; so $\ncl{\pi(S)}_G\subseteq D$.
\end{proof}

Thus $G$ itself already has a nonzero defect containing no torsion, and the
entire question of applying Theorem~\ref{thm:radical} with $K$ as large as
possible reduces to enlarging $\ncl{\pi(S)}$.  Note that no choice of
ingredients makes this automatic: $E/\ncl{P}_E$ is free of rank $2$, so
$\ncl{S}_E$ is never all of $E$, and saturation must be created at a
quotient stage.

\section{Saturation via Hull's theorem}\label{sec:saturation}

Everything in this section is assembled from the published statements
\cite[Definition~1.4, Lemma~3.5, Lemma~5.5, Corollary~5.7, Lemma~5.8,
Theorem~3.12, Theorem~7.1]{Hull16} and \cite[Lemma~7.1]{Osin16}.

\begin{lemma}\label{lem:suitable}
Let $G\in\mathcal{AH}$ be torsion-free and let $N\neq1$ be a normal subgroup.
Fix a generating set $A$ with $\Gamma(G,A)$ hyperbolic and the $G$-action
acylindrical and non-elementary (\cite[Theorem~3.12(4)]{Hull16}).  Then $N$
is suitable with respect to $A$ in the sense of
\cite[Definition~1.4]{Hull16}, and $N$ contains non-commensurable loxodromic
elements $h_1,h_2$ with $E_G(h_i)=\langle h_i\rangle$ such that
$S_0:=\langle h_1,h_2\rangle\leq N$ is again suitable.
\end{lemma}

\begin{proof}
$N$ is nontrivial and torsion-free, hence infinite; being normal it
satisfies $|N^g\cap N|=\infty$ for all $g$, i.e.\ it is $s$-normal, so it
acts non-elementarily by \cite[Lemma~7.1]{Osin16}.  Condition (3) of
suitability ($N$ normalizes no nontrivial finite subgroup) is vacuous in a
torsion-free group, by \cite[Lemma~5.5]{Hull16}.  Now
\cite[Corollary~5.7]{Hull16} applied to the suitable subgroup $N$ yields
$h_1,h_2$ as stated, and \cite[Lemma~5.8]{Hull16} makes
$\langle h_1,h_2\rangle$ suitable.
\end{proof}

\begin{lemma}[defect saturation]\label{lem:saturation}
Let $G\in\mathcal{AH}$ be finitely presented and torsion-free, let
$N\neq1$ be a normal subgroup, and let $F\subseteq G$ be finite.  Then there
is a surjection $q\colon G\twoheadrightarrow Q$ such that
\begin{enumerate}
\item $Q\in\mathcal{AH}$;
\item $Q$ is finitely presented and torsion-free;
\item $q(N)=Q$; indeed $Q=\langle q(h_1),q(h_2)\rangle$ with
  $h_1,h_2\in N$ the elements of Lemma~\textup{\ref{lem:suitable}}, so $Q$
  is $2$-generated;
\item $q|_F$ is injective.
\end{enumerate}
\end{lemma}

\begin{proof}
Let $S_0=\langle h_1,h_2\rangle\leq N$ be as in
Lemma~\ref{lem:suitable}, let $\{t_1,\dots,t_m\}$ be a finite generating set
of $G$, and choose $M$ with $B_A(M)\supseteq F$ (always possible, as Hull
notes).  Apply \cite[Theorem~7.1]{Hull16} to $(G,S_0)$ with these $t_i$ and
this $M$: it returns $q\colon G\twoheadrightarrow Q$ with $Q\in\mathcal{AH}$
(clause (a)), $q|_{B_A(M)}$ injective (clause (b)), and $q(t_i)\in q(S_0)$
for all $i$ (clause (c)).  Then
\[
Q=\langle q(t_1),\dots,q(t_m)\rangle\subseteq q(S_0)\subseteq q(N)\subseteq Q,
\]
so $q(N)=Q=q(S_0)=\langle q(h_1),q(h_2)\rangle$.  Torsion-freeness is clause
(e).  Finite presentation is read off from the proof of Theorem~7.1 exactly
as in \cite{FFF}: the case $m=1$ passes to the quotient by the normal
closure of a single small-cancellation relator
$W=t^{-1}h_1^{m_1}h_2^{l_1}\cdots h_1^{m_n}h_2^{l_n}$ (the set of cyclic
shifts of $W^{\pm1}$ has the same normal closure), and the general case is a
finite induction; so $Q$ is $G$ with $m$ added relators.
\end{proof}

This is the second point we would like an expert to confirm: that
\cite[Corollary~5.7]{Hull16} may be applied to the suitable subgroup $N$ of
Lemma~\ref{lem:suitable} (rather than to $G$ itself) and that the resulting
$S_0\leq N$ feeds into \cite[Theorem~7.1]{Hull16} exactly as Hull's own
proof of \cite[Corollary~7.4]{Hull16} feeds its $S=\langle h_1,h_2\rangle$.
An alternative route with the same output --- re-running Hull's proof of
Corollary~7.4 with the reservoir of target words taken inside a prescribed
infinite normal subgroup of one factor --- is written out in the companion
document (\S\ref{sec:boundary}); the route above has the advantage of using
Theorem~7.1 as a black box applied to the already-constructed group $G$.

\section{The main theorem}\label{sec:main}

\begin{theorem}\label{thm:main}
Let $G$ be the Fournier-Facio group of \S\textup{\ref{sec:ff}} (with
Remark~\textup{\ref{rem:ah}}), and apply Lemma~\textup{\ref{lem:saturation}}
with $N=\ncl{\pi(S)}_G$ and $F=\{1,\pi(s)\}$, obtaining
$q\colon G\twoheadrightarrow Q$.  Then:
\begin{enumerate}
\item $Q$ is $2$-generated, finitely presented, torsion-free,
  acylindrically hyperbolic, and has property \textup{(T)};
\item every nontrivial quotient $L$ of $Q$ --- including $Q$ itself ---
  satisfies $\RadMF(L)=L$; in particular no nontrivial quotient of $Q$ is
  MF, and every homomorphism from $Q$ to an MF group is trivial;
\item every nontrivial quotient of $Q$ is non-sofic.
\end{enumerate}
\end{theorem}

\begin{proof}
(1) All items are Lemma~\ref{lem:saturation} except property \textup{(T)},
which passes to the quotient $Q$ of $G$.

For (2) and (3), let $r\colon Q\twoheadrightarrow L$ be a nontrivial
quotient and write $\rho=r\circ q\circ\pi\colon E\twoheadrightarrow L$.
Since $q(N)=Q$ and images of normal closures are normal closures of images,
\[
L=r(Q)=r\bigl(\ncl{q\pi(S)}_Q\bigr)=\ncl{\rho(S)}_L .
\]
In particular $\rho(S)\neq1$ (else $L=1$), so $\rho|_S$ is injective by
simplicity and $\rho(S)\cong S$.

\emph{The configuration descends.}  Put $\Gamma_L=\rho(P)$,
$t_i^L=\rho(u_i)$, $J_L=(t_1^L)^{-1}\rho(S)\,t_1^L$.  Then
$L=\langle\Gamma_L,t_1^L,t_2^L\rangle$;
$t_i^L\Gamma_L(t_i^L)^{-1}\subseteq\Gamma_L$; $[\Gamma_L,J_L]=1$ and
$t_1^LJ_L(t_1^L)^{-1}=\rho(S)\subseteq\Gamma_L$ (images of the relations in
$G$); $\Gamma_L$ has property \textup{(T)} as a quotient of $P$ and is
infinite since it contains $\rho(S)\cong S$; and $\Gamma_L\cap J_L=1$:
an element $x$ of the intersection lies in $J_L$ and, by
$[\Gamma_L,J_L]=1$, centralizes $J_L$, hence lies in the centre of
$J_L\cong S$, which is trivial.  Finally $L$ has property \textup{(T)} as a
quotient of $Q$.

(2) Take the compression core on $L$ with $\iota$ the inclusion of
$\Gamma_L$, $t=t_1^L$, $c\in J_L\smallsetminus\{1\}$.  By the proof of
Lemma~\ref{lem:defect}, run verbatim in $L$ (using that $\rho(S)\cong S$ is
nonabelian simple with trivial centre), the defect is
$\ncl{\rho(S)}_L=L$.  Apply Theorem~\ref{thm:radical} with $K=L$:
$\RadMF(L)=L$.  Since $L\neq1$, $L$ is not MF; and a nontrivial
homomorphism $Q\to M$ with $M$ MF would make its image a nontrivial MF
quotient... more precisely, its image is a nontrivial quotient $L$ of $Q$
embedded in $M$, and subgroups of MF groups are MF, contradicting (2) for
$L$.

(3) The data of the previous paragraph are exactly the hypotheses of
\cite[Proposition~1.2]{FFF} for $L$, with $J_L\cong S$ finitely generated.
If $L$ were sofic, $S$ would be LEF; but $S$ is finitely presented, infinite
and simple, hence not LEF (\S\ref{sec:obstruction}).
\end{proof}

\begin{corollary}\label{cor:consequences}
$Q$ has no nontrivial finite, amenable, residually finite, linear, sofic, or
MF quotient; $Q$ is perfect; $Q$ has no proper finite-index subgroup; the
profinite and Bohr compactifications of $Q$ are trivial; and every
finite-dimensional unitary representation of $Q$ is trivial.
\end{corollary}

\begin{proof}
Finite, amenable, residually finite, and linear groups (Mal'cev) are sofic
(or directly MF), so a nontrivial quotient of any of these kinds would
contradict Theorem~\ref{thm:main}.  An abelian nontrivial quotient would be
amenable, so $Q$ is perfect.  A proper finite-index subgroup would yield a
nontrivial finite quotient of $Q$ (the action on cosets is nontrivial and
factors through a finite symmetric group, giving a nontrivial finite image).
The compactification and representation statements follow since compact
groups are MF-approximable targets in the relevant sense: a nontrivial
finite-dimensional unitary image would be a nontrivial linear quotient.
\end{proof}

\begin{corollary}\label{cor:marked}
Fix a finite presentation $Q=\langle x,y\mid R\rangle$ and a word $v$ with
$v\neq1$ in $Q$.  The set of $2$-marked groups satisfying the finitely many
relations $R$ and the single inequation $v\neq1$ is a nonempty clopen subset
of the space of $2$-marked groups consisting entirely of groups that are
simultaneously non-sofic and non-MF, each with full MF radical.
Moreover $F_2\twoheadrightarrow Q$ witnesses that neither soficity nor
MF-ness passes to quotients, with a residually finite source; two generators
are optimal, since every $1$-generated group is abelian, hence MF and sofic.
\end{corollary}

\begin{proof}
A marked group satisfying $R$ is a marked quotient of $Q$; the inequation
$v\neq1$ makes it nontrivial; both conditions are clopen.  Apply
Theorem~\ref{thm:main}.
\end{proof}

\begin{remark}
The reduced $C^*$-algebra $C^*_r(Q)$ is stably finite (it has a faithful
trace) and not MF, and no unital $*$-homomorphism
$C^*_{\max}(Q)\to A$ with $A$ an MF algebra can be nontrivial on the
canonical unitaries: it would induce a group homomorphism from $Q$ into
unitaries liftable to matrix models, factoring through the trivial group by
Theorem~\ref{thm:main}.  We deliberately do \emph{not} assert simplicity,
unique trace, or stable rank one of $C^*_r(Q)$ here: those would follow from
the acylindrical-hyperbolicity literature applied to $Q$ (which is
torsion-free with trivial finite radical), but the corresponding inputs are
outside both the formalized development and the scope of this note.
\end{remark}

\section{The certification boundary}\label{sec:boundary}

The accompanying Lean~4 development machine-checks the following, with no
axioms and no literature-transcription premises, in the modules
\begin{center}\ttfamily
Sofic/NormalKazhdanMFRadical,\quad Sofic/TorsionFreeFullMFRadical,\\
Sofic/TorsionFreeFullMFConsequences,\quad Criterion/CriterionAssembly:
\end{center}

\begin{itemize}
\item Theorem~\ref{thm:radical} (\texttt{normalKazhdan\_le\_normMFResidual},
  also in literal Carri\'on--Dadarlat--Eckhardt corona form);
\item the nonsoficity criterion of \S\ref{sec:obstruction}
  (\texttt{not\_isSofic\_of\_not\_isLEF});
\item the defect identification of Lemma~\ref{lem:defect} from the abstract
  Fournier-Facio configuration
  (\texttt{FournierFacioDefectData.simpleSubgroup\_le\_defectNormal});
\item all implications downstream of the saturation step: from
  proof-carrying input data (\texttt{DefectRoutingData}: a surjection with
  saturated defect, two generators, finite presentation, torsion-freeness in
  power form, property \textup{(T)}, one protected element) Lean derives the
  full package of Theorem~\ref{thm:main}(1)--(2) --- including
  $\operatorname{cdeMFResidual}=\top$, failure of the CDE embedding
  property, heredity of the full radical by every quotient, and triviality
  of every homomorphism to an operator-MF group --- and
  Theorem~\ref{thm:main}(3) for the routed quotient, via
  \texttt{NonsoficCriterionData.not\_isSofic} applied to the descended
  criterion data.
\end{itemize}

What is \emph{not} formalized, and is asserted here only on the strength of
the cited literature, is the existence of the input data:
Fournier-Facio's construction (\S\ref{sec:ff}), acylindrical hyperbolicity
and the $\mathcal{AH}_0$ clause (Remark~\ref{rem:ah}), and
Lemmas~\ref{lem:suitable}--\ref{lem:saturation}.  There is deliberately no
Lean theorem of the form ``there exists a group $Q$ with these properties'':
Mathlib has no theory of acylindrically hyperbolic groups or of Hull's
suitable-subgroup small cancellation, and we have chosen to represent that
boundary by structures whose fields are consumed by checked proofs rather
than by axioms.

A companion document in the repository,
\begin{center}
\texttt{docs/TORSION\_FREE\_NORMAL\_GENERATION\_HULL\_QUOTIENT.md},
\end{center}
contains the alternative substitution-based route to
Lemma~\ref{lem:saturation} together with verbatim source quotations of
every Hull and Osin statement used above, extracted from the arXiv PDFs on
August~14, 2026.

\subsection*{The two questions for the expert reader}
\begin{enumerate}
\item Remark~\ref{rem:ah}: is it correct that running
  \cite[\S2]{FFF} with the finitely-generated clause of
  \cite[Corollary~7.4]{Hull16} yields $G\in\mathcal{AH}_0$ with all other
  properties of the construction unchanged?
\item Lemma~\ref{lem:saturation}: is the direct application of
  \cite[Corollary~5.7 and Theorem~7.1]{Hull16} to the suitable subgroup
  $N=\ncl{\pi(S)}_G$ (suitability via \cite[Lemma~7.1]{Osin16} and
  torsion-freeness) correct as stated, including the finite-presentation
  and torsion-preservation readings of the proof of Theorem~7.1?
\end{enumerate}

\begin{thebibliography}{99}

\bibitem[CDE]{CDE}
J.~R. Carri\'on, M.~Dadarlat, and C.~Eckhardt,
\emph{On groups with quasidiagonal $C^*$-algebras},
J. Funct. Anal. \textbf{265} (2013), no.~1, 135--152.

\bibitem[FF]{FFF}
F.~Fournier-Facio,
\emph{A torsion-free non-sofic group},
preprint, 2026, \arxiv{2608.02025}.

\bibitem[Hul]{Hull16}
M.~Hull,
\emph{Small cancellation in acylindrically hyperbolic groups},
Groups Geom. Dyn. \textbf{10} (2016), no.~4, 1077--1119,
\arxiv{1308.4345}.

\bibitem[KT]{KT}
G.~Kun and A.~Thom,
\emph{Nonsofic wreath products of residually finite groups},
preprint, 2026, \arxiv{2608.06222}.

\bibitem[MO]{MO15}
A.~Minasyan and D.~Osin,
\emph{Acylindrical hyperbolicity of groups acting on trees},
Math. Ann. \textbf{362} (2015), no.~3--4, 1055--1105.

\bibitem[OAI]{OAI}
OpenAI,
\emph{A nonsofic group}, 2026.

\bibitem[Os10]{Osin10}
D.~Osin,
\emph{Small cancellations over relatively hyperbolic groups and embedding
theorems},
Ann. of Math. (2) \textbf{172} (2010), no.~1, 1--39,
\arxiv{math/0411039}.

\bibitem[Os16]{Osin16}
D.~Osin,
\emph{Acylindrically hyperbolic groups},
Trans. Amer. Math. Soc. \textbf{368} (2016), no.~2, 851--888,
\arxiv{1304.1246}.

\end{thebibliography}

\end{document}
